\documentclass[10pt,a4paper]{amsart}
\usepackage{amsmath,amssymb,amsthm,graphicx}
% For commutative diagrams we use Xy-pic
\usepackage[all]{xy}

% We use 2cell for 2-commutative diagrams.
\xyoption{2cell}
% Theorem environments.
%
\theoremstyle{plain}
\newtheorem{theorem}[subsection]{Theorem}
\newtheorem{proposition}[subsection]{Proposition}
\newtheorem{lemma}[subsection]{Lemma}

\theoremstyle{definition}
\newtheorem{definition}[subsection]{Definition}
\newtheorem{example}[subsection]{Example}
\newtheorem{exercise}[subsection]{Exercise}
\newtheorem{situation}[subsection]{Situation}

\theoremstyle{remark}
\newtheorem{remark}[subsection]{Remark}
\newtheorem{remarks}[subsection]{Remarks}

\numberwithin{equation}{subsection}

% Macros
%
\def\lim{\mathop{\mathrm{lim}}\nolimits}
\def\colim{\mathop{\mathrm{colim}}\nolimits}
\def\Spec{\mathop{\mathrm{Spec}}}
\def\Hom{\mathop{\mathrm{Hom}}\nolimits}
\def\Ext{\mathop{\mathrm{Ext}}\nolimits}
\def\SheafHom{\mathop{\mathcal{H}\!\mathit{om}}\nolimits}
\def\SheafExt{\mathop{\mathcal{E}\!\mathit{xt}}\nolimits}
\def\Sch{\mathit{Sch}}
\def\Mor{\mathop{\mathrm{Mor}}\nolimits}
\def\Ob{\mathop{\mathrm{Ob}}\nolimits}
\def\Sh{\mathop{\mathit{Sh}}\nolimits}
\def\NL{\mathop{N\!L}\nolimits}
\def\CH{\mathop{\mathrm{CH}}\nolimits}
\def\proetale{{pro\text{-}\acute{e}tale}}
\def\etale{{\acute{e}tale}}
\def\QCoh{\mathit{QCoh}}
\def\Ker{\mathop{\mathrm{Ker}}}
\def\Im{\mathop{\mathrm{Im}}}
\def\Coker{\mathop{\mathrm{Coker}}}
\def\Coim{\mathop{\mathrm{Coim}}}

% Boxtimes
%
\DeclareMathSymbol{\boxtimes}{\mathbin}{AMSa}{"02}

%
% Macros for moduli stacks/spaces
%
\def\QCohstack{\mathcal{QC}\!\mathit{oh}}
\def\Cohstack{\mathcal{C}\!\mathit{oh}}
\def\Spacesstack{\mathcal{S}\!\mathit{paces}}
\def\Quotfunctor{\mathrm{Quot}}
\def\Hilbfunctor{\mathrm{Hilb}}
\def\Curvesstack{\mathcal{C}\!\mathit{urves}}
\def\Polarizedstack{\mathcal{P}\!\mathit{olarized}}
\def\Complexesstack{\mathcal{C}\!\mathit{omplexes}}
% \Pic is the operator that assigns to X its picard group, usage \Pic(X)
% \Picardstack_{X/B} denotes the Picard stack of X over B
% \Picardfunctor_{X/B} denotes the Picard functor of X over B
\def\Pic{\mathop{\mathrm{Pic}}\nolimits}
\def\Picardstack{\mathcal{P}\!\mathit{ic}}
\def\Picardfunctor{\mathrm{Pic}}
\def\Deformationcategory{\mathcal{D}\!\mathit{ef}}
\usepackage{cleveref}
\crefname{originalreference}{original reference}{original references}
\Crefname{originalreference}{Original reference}{Original references}
\makeatletter
\newlabel{lemma-finite-cd}{{0F10}{source}{}{sourcepointer-1}{}}
\newlabel{lemma-finite-cd@cref}{{[originalreference][0][]0F10}{[1][0][]source}{}{}{}}
\makeatother
\crefname{originalreference}{original reference}{original references}
\Crefname{originalreference}{Original reference}{Original references}
\makeatletter
\newlabel{lemma-finite-cd-mod-n-direct-sums}{{0F11}{source}{}{sourcepointer-2}{}}
\newlabel{lemma-finite-cd-mod-n-direct-sums@cref}{{[originalreference][0][]0F11}{[1][0][]source}{}{}{}}
\makeatother
\crefname{originalreference}{original reference}{original references}
\Crefname{originalreference}{Original reference}{Original references}
\makeatletter
\newlabel{lemma-punctual-base-change-upgrade}{{0F1J}{source}{}{sourcepointer-3}{}}
\newlabel{lemma-punctual-base-change-upgrade@cref}{{[originalreference][0][]0F1J}{[1][0][]source}{}{}{}}
\makeatother
\title{The punctual base change isomorphism for unbounded complexes of Z/nZ-modules over a field}
\author{The Stacks Project Authors}
\date{stacks-project:etale-cohomology.tex:local-snapshot-20261011; GNU FDL 1.2 or later}
\begin{document}
\maketitle
\noindent\textit{Editorial note.} All delivered mathematical text is the author's own, copied from the pinned Stacks Project chapter file etale-cohomology.tex; no statement or proof step was rewritten, reordered or removed.\par\medskip
\noindent\textit{Editorial note.} The counted proof is the author's complete argument: the proof is written out in full: the list of commutation-with-direct-sums inputs, the homotopy colimit and truncation reductions for E and F, the passage to a bounded complex of flat representatives of F, and the final appeal to the bounded punctual base change statement. The author's own remark that the reader may skip the proof is retained in the delivered text and is a reading suggestion, not an omission: no step is left out.\par\medskip
\noindent\textit{Editorial note.} The label lemma-punctual-base-change-upgrade-unbounded names the counted statement, and the proof environment immediately following it in the source is delivered verbatim as its proof.\par\medskip
\noindent\textit{Editorial note.} The counted statement is carried by the author's own lemma environment in the chapter source.\par\medskip
\noindent\textit{Editorial note.} The statement and proof are copied from the pinned chapter file; the author's sentence suggesting the reader skip the proof is part of the delivered text, and the proof is nonetheless complete.\par\medskip
\noindent\textit{Editorial note.} The preamble is the chapter's own native preamble excerpt already used for this source: the Xy-pic setup (native preamble lines 14-18) and the native theorem declarations with the mathematical macros (native preamble lines 154-217).\par\medskip
\noindent\textit{Allowed standard prerequisites (editorial).}\begin{itemize}
\item The bounded punctual base change statement. Same chapter, the punctual base change lemma in its bounded form for complexes of Z/nZ-modules, invoked at the end of the proof. Supplies the case to which the unbounded statement is reduced.
\item Commutation of the relevant functors with direct sums. Same chapter, the two lemmas stating that the derived tensor product, the pullback, and the pushforwards along the finite type morphisms commute with direct sums. Supplies the passage from arbitrary to bounded above complexes, and through homotopy colimits from arbitrary objects of the derived category.
\item Truncations and flat resolutions of complexes bounded above. Derived Categories and same chapter, the stupid truncations and the existence of bounded above flat resolutions used to reduce F. Supplies the reduction to a bounded complex for E and to a bounded complex of flat sheaves for F.
\end{itemize}\medskip
\ifdefined\hypertarget\hypertarget{sourcepointer-1}{}\fi
\noindent\textit{Source reference.} Original reference 0F10: sources/stacks-etale-cohomology/pinned-original-source.tex, source lines 20041--20041. Prerequisite pointer invoked at proof line(s) 20908; it names a standard fact of the same chapter that the counted argument uses, and every step of the counted proof built on it is retained verbatim in the delivered body.\par\medskip
\ifdefined\hypertarget\hypertarget{sourcepointer-2}{}\fi
\noindent\textit{Source reference.} Original reference 0F11: sources/stacks-etale-cohomology/pinned-original-source.tex, source lines 20074--20074. Prerequisite pointer invoked at proof line(s) 20909; it names a standard fact of the same chapter that the counted argument uses, and every step of the counted proof built on it is retained verbatim in the delivered body.\par\medskip
\ifdefined\hypertarget\hypertarget{sourcepointer-3}{}\fi
\noindent\textit{Source reference.} Original reference 0F1J: sources/stacks-etale-cohomology/pinned-original-source.tex, source lines 20706--20706. Prerequisite pointer invoked at proof line(s) 20902; it names a standard fact of the same chapter that the counted argument uses, and every step of the counted proof built on it is retained verbatim in the delivered body.\par\medskip
\begin{lemma}
\label{lemma-punctual-base-change-upgrade-unbounded}
Let $K$ be a field. Let $n \geq 1$ be invertible in $K$.
Consider a commutative diagram
$$
\xymatrix{
X \ar[d] & X' \ar[l]^p \ar[d]_{f'} & Y \ar[l]^h \ar[d]^e \\
\Spec(K) & S' \ar[l] & T \ar[l]_g
}
$$
of schemes of finite type over $K$ with
$X' = X \times_{\Spec(K)} S'$ and $Y = X' \times_{S'} T$.
The canonical map
$$
p^{-1}E \otimes_{\mathbf{Z}/n\mathbf{Z}}^\mathbf{L} (f')^{-1}Rg_*F
\longrightarrow
Rh_*(h^{-1}p^{-1}E \otimes_{\mathbf{Z}/n\mathbf{Z}}^\mathbf{L} e^{-1}F)
$$
is an isomorphism for $E$ in $D(X_\etale, \mathbf{Z}/n\mathbf{Z})$
and $F$ in $D(T_\etale, \mathbf{Z}/n\mathbf{Z})$.
\end{lemma}

\begin{proof}
We will reduce this to
Lemma \ref{lemma-punctual-base-change-upgrade}
using that our functors commute with direct sums.
We suggest the reader skip the proof.
Recall that derived tensor product commutes with
direct sums. Recall that (derived) pullback commutes with direct sums.
Recall that $Rh_*$ and $Rg_*$ commute with direct sums, see
Lemmas \ref{lemma-finite-cd} and
\ref{lemma-finite-cd-mod-n-direct-sums}
(this is where we use our schemes are of finite type
over $K$).

\medskip\noindent
To finish the proof we can argue as follows.
First we write $E = \text{hocolim} \tau_{\leq N} E$.
Since our functors commute with direct sums, they commute
with homotopy colimits. Hence it suffices to prove
the lemma for $E$ bounded above. Similarly for $F$ we
may assume $F$ is bounded above.
Then we can represent $E$ by a bounded above complex
$\mathcal{E}^\bullet$ of sheaves of $\mathbf{Z}/n\mathbf{Z}$-modules.
Then
$$
\mathcal{E}^\bullet = \colim \sigma_{\geq -N}\mathcal{E}^\bullet
$$
(stupid truncations).
Thus we may assume $\mathcal{E}^\bullet$ is a bounded
complex of sheaves of $\mathbf{Z}/n\mathbf{Z}$-modules.
For $F$ we choose a bounded above complex of
flat(!) sheaves of $\mathbf{Z}/n\mathbf{Z}$-modules.
Then we reduce to the case where $F$ is represented
by a bounded complex of flat sheaves of $\mathbf{Z}/n\mathbf{Z}$-modules.
At this point Lemma \ref{lemma-punctual-base-change-upgrade}
kicks in and we conclude.
\end{proof}

\end{document}
